\chapter{Discussion}
\label{chap:discussion}

\section{Implementation}

The introduction of a filter on the velocity was found to be necessary during testing. The implementation uses an exponential moving average which can be tuned using a scalar $\alpha$. It is important to note that this is trade-off between noise suppression and phase delay. It was vital in order to limit jittering, while the phase delay did not seem to impact performance to a point of instability even with an $\alpha = 0$.

The implementation of the controllers on the physical system does come with some challenges and limitations. One of the early important factors to consider was the minimum cycle time of the PLC. The default cycle time is set at 10 ms, which quickly showed itself as an issue. With the first attempt at getting a controller to work, a preliminary PID controller was tested at the default cycle time. Unsurprisingly it did not stabilize. After some trial and error, changing parameters, gains and the like, to no avail, it was found that the cycle time was the crucial issue. Lowering it first to 2 ms, got the pendulum somewhat stable. Later the cycle time was pushed to the hardware limit of 1 ms, with a short tolerance for cycle time overshoot. This allows the PLC to continue running its primary loop even if the computation time exceeds the cycle time, with a margin of 10 ms. If the total delay ever got to 10 ms, the loop would stop running. At no point during the testing phase did the PLC shut off due to this computation delay. 

While it was not achievable to get the pendulum stable at a 100 Hz sampling frequency during preliminary testing, it should theoretically be possible. One benefit to using a slower sampling rate is the reduction of noise impact. With a very fast sampling rate, any slight changes and inconsistencies in the data will impact the calculations substantially. A slower rate will act as a natural filter, smoothing some of the noise, but also increasing the delay of information.

\section{Data Collection}
One of the major hardware limitations of the PLC is its storage capacity. This means it is not viable to collect the data directly on the PLC to be moved later. Thus the data must be moved from the PLC to another device continuously, which is why an OPC UA system was implemented. The speed of the OPC UA communication is somewhat limited, and paired with python, slightly fluctuating. This shows itself in the test data as the data points are not collected at a fixed rate. Though the PLC sampled information every millisecond, the collected samples have 3-4 ms between them. In order to counteract this uncertainty, the samples are sent as the last thing during a cycle and a step count from the PLC is sent along with it. This means the state of system should match the step count and should be consistent in each data sample, allowing the data to be synchronized. While this solution is sub-optimal, it was determined to be acceptable for the testing required in this project.